\section{Patrón arquitectónico}

La arquitectura adopta un flujo unidireccional inspirado en Redux, ajustado al ciclo de vida de Unity. Los datos actuales residen en un \emph{store}; cualquier intención de cambio se expresa como una acción; un reductor calcula el siguiente estado; finalmente, un adaptador aplica al runtime los efectos necesarios.

\subsection{Componentes y responsabilidades}

\begin{table}[H]
\centering
\caption{Responsabilidades de los componentes del patrón}
\begin{tabularx}{\textwidth}{>{\bfseries}p{3.1cm}X X}
\toprule
Componente & Responsabilidad & Restricción principal \\
\midrule
StateData & Contener una instantánea serializable del dominio. & No controla objetos de escena ni ejecuta efectos. \\
Action & Nombrar una intención: mover, cambiar rol, configurar colisión. & Debe ser explícita y tener un tipo inequívoco. \\
Payload & Transportar los valores necesarios para la acción. & Solo datos; sin acceso al store. \\
Reducer & Calcular el siguiente estado desde el estado actual y la acción. & Debe ser determinista y no modificar la escena. \\
Store & Custodiar el estado, ejecutar el reductor y publicar cambios. & Es el único punto de escritura del estado. \\
Runtime Applier & Traducir el estado a componentes Unity. & Escucha al store; no decide reglas de dominio. \\
Controlador & Lee entrada o eventos y despacha acciones. & No escribe directamente el StateData. \\
\bottomrule
\end{tabularx}
\end{table}

\subsection{Flujo unidireccional}

\begin{figure}[H]
  \centering
  \includegraphics[width=0.88\textwidth]{\figurepath/flujo-datos.png}
  \caption{Flujo conceptual desde el estado actual hasta la actualización del store y la reconciliación física.}
  \label{fig:flujo-datos}
\end{figure}

El flujo esperado es:

\begin{enumerate}
  \item un controlador captura una intención del usuario o un evento del juego;
  \item calcula o solicita un valor permitido, por ejemplo una posición sin penetrar colisiones;
  \item crea una \code{Action} con su \code{Payload};
  \item el \code{Store.Dispatch} entrega la acción al \code{Reducer};
  \item el reductor valida reglas simples y devuelve una nueva instantánea;
  \item si el estado cambió, el store publica \code{StateChanged};
  \item los adaptadores actualizan \code{Rigidbody2D}, \code{BoxCollider2D}, renderizadores u otros componentes;
  \item cuando Unity altera el valor final por física, se despacha una acción de reconciliación.
\end{enumerate}

\subsection{Propiedades del diseño}

\begin{itemize}
  \item \textbf{Trazabilidad:} cada mutación se identifica mediante un tipo de acción.
  \item \textbf{Repetibilidad:} el mismo estado y la misma acción producen el mismo resultado.
  \item \textbf{Desacoplamiento:} la lógica de datos no necesita referencias a objetos de escena.
  \item \textbf{Pruebas simples:} reductores y constructores se prueban sin cargar una escena.
  \item \textbf{Integración segura:} los efectos de Unity quedan concentrados en controladores y adaptadores.
\end{itemize}

\subsection{Inmutabilidad práctica}

En el código actual, \code{PlayerStateData} y \code{GameSessionStateData} son estructuras serializables con campos privados y propiedades de solo lectura. Los métodos \code{With...} construyen una nueva instantánea en vez de exponer setters. Este enfoque debe mantenerse para los estados planeados.

\begin{lstlisting}[caption={Forma recomendada de producir una nueva instantánea},label={lst:with-pattern}]
public WorldStateData WithRoom(string value) =>
    new WorldStateData(id, name, value, mainTheme, dimensions);
\end{lstlisting}

El uso de métodos \code{set.X} y \code{get.X} en los diagramas representa la intención funcional. En C\# se recomienda expresarla mediante fábricas de acciones, propiedades de solo lectura y métodos \code{With...}; no mediante setters públicos que permitan saltarse el store.
